\documentclass[11pt]{article}
\usepackage[margin=1in]{geometry}
\usepackage{amsmath, amssymb, amsthm}
\usepackage{hyperref}

\title{Options Pricing Theory:\\Black-Scholes, Binomial Trees, Monte Carlo, and the Greeks}
\author{Abhi Wadhwa}
\date{}

\begin{document}
\maketitle

\begin{abstract}
These notes cover the mathematical foundations of options pricing: the Black-Scholes model, the Cox-Ross-Rubinstein binomial tree, Monte Carlo simulation, and the Greeks. I try to explain not just the formulas but why they work---what's really going on with risk-neutral pricing, why the drift disappears, and how the different methods connect. I also cover put-call parity and implied volatility, which are essential tools for any practitioner.
\end{abstract}

\section{Geometric Brownian Motion}

Everything starts with a model for the stock price. The standard assumption is \textbf{Geometric Brownian Motion} (GBM):
\[
dS = \mu S \, dt + \sigma S \, dW
\]
where $\mu$ is the drift (expected return), $\sigma$ is the volatility, and $W$ is a standard Brownian motion. The key feature of GBM is that returns are log-normal: $\ln(S_T/S_0) \sim \mathcal{N}\big((\mu - \sigma^2/2)T,\; \sigma^2 T\big)$.

Why GBM? Because it guarantees $S > 0$ (stock prices can't go negative), and it captures the empirical observation that percentage returns are roughly normally distributed. It's not perfect---real returns have fat tails and volatility clustering---but it's a remarkably good first-order model.

\section{The Black-Scholes Formula}

\subsection{Risk-Neutral Pricing}

Here's the part that confused me for the longest time: \textbf{the drift $\mu$ disappears from the option price}. This seems wrong---surely an option on a stock that's expected to go up is worth more?

The resolution is risk-neutral pricing. The idea is that you can perfectly replicate the option payoff by dynamically trading the stock and a bond. Since the replicating portfolio and the option have the same payoff, they must have the same price (no arbitrage). When you work through the replication argument, the drift cancels out because both the option and the stock are driven by the same Brownian motion.

The result is that we can price options as if we're in a ``risk-neutral world'' where the stock grows at the risk-free rate $r$ instead of $\mu$:
\[
dS = rS \, dt + \sigma S \, dW^{\mathbb{Q}}
\]

\subsection{Derivation Sketch via Ito's Lemma}

Let $V(S,t)$ be the option price as a function of stock price and time. By Ito's lemma:
\[
dV = \frac{\partial V}{\partial t} dt + \frac{\partial V}{\partial S} dS + \frac{1}{2} \frac{\partial^2 V}{\partial S^2} (dS)^2
\]

Since $(dS)^2 = \sigma^2 S^2 dt$ (the $dt \cdot dt$ and $dt \cdot dW$ terms vanish), this becomes:
\[
dV = \left(\frac{\partial V}{\partial t} + \mu S \frac{\partial V}{\partial S} + \frac{1}{2}\sigma^2 S^2 \frac{\partial^2 V}{\partial S^2}\right) dt + \sigma S \frac{\partial V}{\partial S} dW
\]

Now construct a portfolio $\Pi = V - \Delta S$ where $\Delta = \partial V / \partial S$. The $dW$ terms cancel (that's the whole point of delta hedging), leaving a deterministic portfolio that must earn the risk-free rate. Setting $d\Pi = r\Pi \, dt$ gives the \textbf{Black-Scholes PDE}:
\[
\boxed{\frac{\partial V}{\partial t} + \frac{1}{2}\sigma^2 S^2 \frac{\partial^2 V}{\partial S^2} + rS\frac{\partial V}{\partial S} - rV = 0}
\]

Notice: $\mu$ is gone. The drift doesn't matter for option pricing.

\subsection{The Formulas}

Solving the BS PDE with the boundary condition $V(S,T) = \max(S-K, 0)$ for a call gives:
\[
\boxed{C = S \cdot N(d_1) - K e^{-rT} \cdot N(d_2)}
\]
where
\[
d_1 = \frac{\ln(S/K) + (r + \sigma^2/2)T}{\sigma\sqrt{T}}, \qquad d_2 = d_1 - \sigma\sqrt{T}
\]
and $N(\cdot)$ is the standard normal CDF.

For a put:
\[
P = K e^{-rT} \cdot N(-d_2) - S \cdot N(-d_1)
\]

The intuition for the call formula: $N(d_2)$ is (roughly) the risk-neutral probability that the option expires in the money, so $K e^{-rT} N(d_2)$ is the PV of the expected strike payment. $S \cdot N(d_1)$ is the expected present value of receiving the stock, adjusted for the correlation between the stock price and the event of exercise.

\section{Put-Call Parity}

Before moving on, here's a result that's independent of any model:
\[
\boxed{C - P = S - K e^{-rT}}
\]

This follows from a simple no-arbitrage argument: a long call and short put with the same strike and expiry is equivalent to a forward contract. If this relationship is violated, you can construct a riskless arbitrage.

I use put-call parity as a sanity check constantly. If my call and put prices don't satisfy this exactly, something is wrong with my code.

\section{The Binomial Model (CRR)}

The Cox-Ross-Rubinstein (1979) binomial tree is the discrete-time version of Black-Scholes, and honestly I think it's the best way to build intuition for risk-neutral pricing.

\subsection{One-Step Model}

At each time step $\Delta t$, the stock either goes up by factor $u$ or down by factor $d$:
\[
u = e^{\sigma\sqrt{\Delta t}}, \qquad d = \frac{1}{u} = e^{-\sigma\sqrt{\Delta t}}
\]

The risk-neutral probability of an up move is:
\[
p = \frac{e^{r\Delta t} - d}{u - d}
\]

The trick here is that $p$ is chosen so that the expected return under this probability is exactly $r$:
\[
p \cdot u + (1-p) \cdot d = e^{r\Delta t}
\]

The option value at any node is then the discounted expected value under $p$:
\[
V = e^{-r\Delta t}\big[p \cdot V_{\text{up}} + (1-p) \cdot V_{\text{down}}\big]
\]

\subsection{Multi-Step Tree}

Build a recombining tree with $n$ steps. At expiry, the option value at each terminal node is the intrinsic value: $\max(S \cdot u^j d^{n-j} - K, 0)$ for a call. Then work backwards through the tree using the one-step formula.

For \textbf{American options}, at each node we also check whether early exercise is optimal:
\[
V_{\text{node}} = \max\!\big(\text{exercise value},\; e^{-r\Delta t}[p \cdot V_{\text{up}} + (1-p) \cdot V_{\text{down}}]\big)
\]

As $n \to \infty$, the binomial price converges to the Black-Scholes price. This is a beautiful result and it's one of those things where you can actually watch it happen by running the code with increasing $n$.

\section{Monte Carlo Simulation}

Monte Carlo is the most flexible pricing method. You simulate many stock price paths under the risk-neutral measure and average the discounted payoffs.

Each path is simulated as:
\[
S(t + \Delta t) = S(t) \cdot \exp\!\left[\left(r - \frac{\sigma^2}{2}\right)\Delta t + \sigma\sqrt{\Delta t} \cdot Z\right]
\]
where $Z \sim \mathcal{N}(0,1)$. The option price is:
\[
V = e^{-rT} \cdot \frac{1}{N}\sum_{i=1}^{N} \text{payoff}(S_T^{(i)})
\]

For a European call, $\text{payoff}(S_T) = \max(S_T - K, 0)$.

Monte Carlo converges as $O(1/\sqrt{N})$---you need 100x more paths to get 10x more precision, which is slow. But it handles exotic payoffs (Asian options, barrier options, lookbacks) trivially, whereas binomial trees and closed-form solutions often can't.

After playing with this for a while, I found that \textbf{antithetic variates} (for each $Z$, also simulate $-Z$) roughly halve the variance for free. It's one of those easy wins that should always be used.

\section{The Greeks}

The Greeks measure how the option price changes with respect to each input parameter. They're essential for risk management and hedging.

\subsection{Delta: $\Delta = \partial V / \partial S$}

For a BS call, $\Delta = N(d_1)$. Delta ranges from 0 (deep OTM) to 1 (deep ITM). It tells you how many shares to hold to hedge the option.

\textbf{Key fact}: $\Delta_{\text{call}} - \Delta_{\text{put}} = 1$. This follows directly from put-call parity.

\subsection{Gamma: $\Gamma = \partial^2 V / \partial S^2$}

Gamma measures how fast delta changes. High gamma means you need to rebalance your hedge frequently. Gamma is always positive for long options and is highest near the strike when the option is close to expiry.
\[
\Gamma = \frac{n(d_1)}{S\sigma\sqrt{T}}
\]
where $n(\cdot)$ is the standard normal PDF.

\subsection{Theta: $\Theta = \partial V / \partial t$}

Time decay. Options lose value as they approach expiry (all else equal). For a call:
\[
\Theta = -\frac{S \cdot n(d_1) \cdot \sigma}{2\sqrt{T}} - rKe^{-rT}N(d_2)
\]

There's a beautiful relationship between the Greeks that comes directly from the BS PDE:
\[
\Theta + \frac{1}{2}\sigma^2 S^2 \Gamma + rS\Delta = rV
\]

This is just the PDE rewritten in terms of Greeks. It means theta and gamma are linked: an option with high gamma (lots of convexity) must have high theta (lots of time decay) to compensate.

\subsection{Vega: $\mathcal{V} = \partial V / \partial \sigma$}

Sensitivity to volatility. Always positive for long options:
\[
\mathcal{V} = S \cdot n(d_1) \cdot \sqrt{T}
\]

\subsection{Rho: $\rho = \partial V / \partial r$}

Sensitivity to the risk-free rate. Usually small:
\[
\rho_{\text{call}} = KTe^{-rT}N(d_2)
\]

\section{Implied Volatility}

In practice, we observe market prices and want to extract the ``implied'' volatility---the $\sigma$ that, when plugged into the BS formula, reproduces the observed price. This is solving:
\[
\text{BS}(S, K, T, r, \sigma_{\text{impl}}) = V_{\text{market}}
\]

Since BS is monotonically increasing in $\sigma$ (vega is always positive), this equation has a unique solution and can be found numerically. Newton-Raphson is fast because we have analytic vega:
\[
\sigma_{n+1} = \sigma_n - \frac{\text{BS}(\sigma_n) - V_{\text{market}}}{\text{vega}(\sigma_n)}
\]

Brent's method is slower but more robust (doesn't require derivatives and always converges given a valid bracket).

Implied volatility is arguably the most important quantity in options markets. Traders think in terms of vol, not price. The fact that implied vol varies across strikes (the ``volatility smile'') tells us that the market doesn't really believe in the log-normal assumption---but Black-Scholes remains the lingua franca.

\begin{thebibliography}{9}
\bibitem{black1973} Black, F., Scholes, M. (1973). ``The Pricing of Options and Corporate Liabilities.'' \textit{Journal of Political Economy}, 81(3), 637--654.

\bibitem{cox1979} Cox, J.C., Ross, S.A., Rubinstein, M. (1979). ``Option Pricing: A Simplified Approach.'' \textit{Journal of Financial Economics}, 7(3), 229--263.

\bibitem{hull2018} Hull, J.C. (2018). \textit{Options, Futures, and Other Derivatives}. 10th edition, Pearson.

\bibitem{merton1973} Merton, R.C. (1973). ``Theory of Rational Option Pricing.'' \textit{Bell Journal of Economics and Management Science}, 4(1), 141--183.

\bibitem{glasserman2004} Glasserman, P. (2004). \textit{Monte Carlo Methods in Financial Engineering}. Springer.
\end{thebibliography}

\end{document}
